\documentclass[11pt]{article}
\usepackage[a4paper,margin=21mm]{geometry}
\usepackage{fontspec}
\setmainfont{Latin Modern Roman}
\newfontfamily\CJKfont{Noto Serif CJK SC}
\XeTeXlinebreaklocale "zh"
\XeTeXlinebreakskip=0pt plus 1pt
\usepackage{amsmath,amssymb,amsthm,amscd,graphicx,color}
\usepackage{xurl}
\usepackage[unicode,hidelinks]{hyperref}
\allowdisplaybreaks
\setlength\emergencystretch{3em}
\numberwithin{equation}{section}

\usepackage{mathtools}

\usepackage[all]{xy}

\numberwithin{equation}{section}

\newtheorem{Theorem}{Theorem}[section]
\newtheorem*{Theorem*}{Theorem}
\newtheorem{Corollary}[Theorem]{Corollary}
\newtheorem{Lemma}[Theorem]{Lemma}
\newtheorem{Proposition}[Theorem]{Proposition}
\newtheorem{Construction}[Theorem]{Construction}

 { \theoremstyle{definition}
\newtheorem{Definition}[Theorem]{Definition}
\newtheorem{Note}[Theorem]{Note}
\newtheorem{Example}[Theorem]{Example}
\newtheorem{Remark}[Theorem]{Remark} }


\def\CC{\mathbb{C}}
\def\DD{\mathbb{D}}
\def\LL{\mathbb{L}}
\def\PP{\mathbb{P}}
\def\RR{\mathbb{R}}
\def\ZZ{\mathbb{Z}}

%%%%%%%% mathcal %%%%%%%%

\def\cW{\mathcal{W}}
\def\calP{\mathcal{P}}
\def\calL{\mathcal{L}}
\def\cC{\mathcal{C}}
\def\cO{\mathcal{O}}
\def\cF{\mathcal{F}}
\def\cE{\mathcal{E}}

%%%%%%%% frak %%%%%%%%

\newcommand\frX{\mathfrak{X}}
\newcommand\frs{\mathfrak{s}}

\newcommand{\Coh}{\textup{Coh}}
\newcommand{\Ind}{\textup{Ind}}
\newcommand\Loc{\textup{Loc}}
\newcommand{\QCoh}{\textup{QCoh}}
\newcommand\Spec{\textup{Spec}}
\newcommand\Hom{\textup{Hom}}

%%%%%%%% Lie groups and algebras %%%%%%%
\newcommand\SL{\textup{SL}}

\newcommand\nc{\newcommand}
\nc\oX{\overline{X}}
\nc\oD{\overline{D}}
\nc\oT{\overline{T}}
\nc\bsig{\overline{\sigma}}
\nc\Tors{\on{Tors}}

\nc\on{\operatorname}
\nc\ol{\overline}
\nc\ul{\underline}
\nc\Bl{\on{Bl}}
\nc\Conv{\on{Conv}}
\nc\Int{\on{Int}}
\nc\Pc{\mathring{P}}

\nc\Cone{\on{Cone}}
\nc\colim{\varinjlim}
\nc\Sh{\mathrm{Sh}}
\nc\wsh{\Sh^w}
\nc\wmsh{\mu\Sh^w}
\nc\mmod{\on{-mod}}
\nc\ovi{\overline{i}}


\makeatletter\newlabel{subsec:future}{{1.3}{}{Original source locator}{}{}}\makeatother
\makeatletter\newlabel{ex:basic-surgery-highd}{{4.5}{}{Original source locator}{}{}}\makeatother
\makeatletter\newlabel{ex:fund-calc}{{4.11}{}{Original source locator}{}{}}\makeatother
\begin{document}
\begin{center}\Large Coherent-sheaf dg categories of blowups with the transformed divisor deleted are the stated homotopy pushouts for a smooth codimension-two center in a divisor\end{center}
\noindent\textbf{Stable ID:} 1249\quad\textbf{Author:} Benjamin Gammage; Ian Le\par
\noindent\textbf{Work:} Mirror Symmetry for Truncated Cluster Varieties\par
\noindent\textbf{Source:} \url{https://sigma-journal.com/2022/055/}\par
\noindent\textbf{Version:} Official SIGMA 18 (2022) article 055, DOI 10.3842/SIGMA.2022.055; journal PDF and native TeX acquired 2026-09-30, exact bytes pinned by SHA256\par
\noindent\textbf{Rights:} CC BY 4.0. Authors retain copyright. \url{https://creativecommons.org/licenses/by/4.0/}. Reuse and typesetting adaptation; original language and mathematical content preserved.\par
\noindent\textbf{Difficulty (editorial):} {\CJKfont H2: The proof constructs two geometric functors from the center and ambient variety and identifies their restrictions only after deleting the transformed divisor. It then uses the blowup semiorthogonal decomposition to prove essential surjectivity and identifies the quotient by divisor-supported objects with the claimed homotopy coequalizer. The new geometric localization-to-pushout comparison is substantially more than applying the prior blowup decomposition alone.}\par
\medskip\noindent\textbf{Editorial boundary note (not author text):} Selected complete current Theorem4.13 geometric homotopy pushout, for smooth variety or DM stack X, divisor D, and smooth H contained in D of codimension2 in X. Characteristic0 stable-infinity/dg conventions, all derived functors, blowup-minus-strict-transform space and i-star/k-pushforward pushout are retained. Full author FH/FX construction, agreement after restriction, diagram, semiorthogonal generation, localization and coproduct/coequalizer argument are unchanged. Orlov blowup decomposition and standard divisor localization are established author prerequisites. The printed revised-version notice is preserved exactly; removed old Section5 is not part of this selected current proof. HMS4.12 and helper mirror/symplectic examples receive no counts. Native XY diagram and every functor/sign/twist are editable and literal.\par
\noindent\textbf{External original reference subsec:future:} Current revised-version future questions native162–177, exact original revised notice119 retained; selected Theorem4.13 does not use removed Section5\par
\noindent\textbf{External original reference ex:basic-surgery-highd:} Unselected Weinstein-attachment illustration native643–685; not a selected geometric pushout premise\par
\noindent\textbf{External original reference ex:fund-calc:} Unselected local mirror illustration native715–745; selected theorem is stated for general smooth variety/DM stack\par
\medskip\hrule\medskip\noindent\textbf{Author text begins below.}\par

% Original source lines 119--119
\textit{Note on revised version: an earlier version of this paper contained an error in Section~{\rm 5}, which invalidated the proofs of some results in Sections~{\rm 5.2--5.3}; we have now removed that section of the paper, and we now discuss our expectations on the relation to toric geometry at the end of Section~{\rm \ref{subsec:future}}.}


% Original source lines 188--189
\subsection*{Categorical conventions}
Throughout this paper, we work with stable $\infty$-categories over a field $k$ of characteristic 0, which we model as dg categories (so that, for instance, $\Coh(X)$ always refers to the dg category of coherent sheaves on $X$). All functors are derived, and the word ``limits/colimit'' in a diagram of categories is always used to refer to a homotopy limit/colimit.

\setcounter{section}{4}\setcounter{Theorem}{12}\setcounter{equation}{3}
% Original source lines 767--812
We will deduce this theorem from the following more general result, which we can understand as a generalization of the calculation described in Example~\ref{ex:fund-calc}:

\begin{Theorem}\label{thm:main-general}
	Let $X$ be a smooth variety or DM stack, $D\xhookrightarrow{k} X$ a divisor, and $H\xhookrightarrow{i}D$ a~smooth subvariety of codimension $2$ in $X$. Let $\frX = \Bl_HX\setminus \widetilde{D}$ and let $\cC$ denote the
homotopy pushout of dg categories
\begin{equation}\label{eq:pushout}
\colim\bigl(\Coh(H)\stackrel{i^*}\longleftarrow\Coh(D)\stackrel{k_*}\longrightarrow\Coh(X)\bigr).
\end{equation}
Then there is an equivalence
\[
\cC\cong\Coh(\frX)
\] between $\cC$ and the category of coherent sheaves on $\frX$.
\end{Theorem}
\begin{proof}
We begin by writing down the functor $\cC\to \Coh(\frX)$. From the pushout presentation of $\cC$, such a functor may be specified by the data of functors $F_H\colon \Coh(H)\to \Coh(\frX)$ and $F_X\colon \Coh(X)\to\Coh(\frX)$, together with an equivalence between the pullbacks of $F_H$ and $F_X$ to $\Coh(D)$ along the maps in the pushout \eqref{eq:pushout}.

Consider the following diagram of all the spaces involved in our story, where $E$ is the exceptional divisor of the blowup, the vertical maps are the projections, and the horizontal maps are the inclusions:
\[
\xymatrix{
\frX\ar@{^{(}->}[r]^-j & \Bl_HX\ar[d]_p & &E\ar@{_{(}->}[ll]_-{\ovi}\ar[d]^q\\
& X & D\ar@{_{(}->}[l]^k & H.\ar@{_{(}->}[l]^-i
}
\]
The functors $F_H$ and $F_X$ will be defined as the compositions
\begin{gather*}
F_H = j^*\ovi_* q^*, \qquad F_X = j^*p^*.
\end{gather*}
The restriction to $D$ of $F_X$ is evidently given by $j^*p^*k_*$, and we must produce an isomorphism between this functor and the other restriction $j^*\ovi_*q^*i^*$. Note that before applying final composition with $j^*$, these two functors differ, since $p^*k_*$ takes sheaves on $D$ to their pullback to the blowup, whereas $\ovi_* q^*i^*$ takes sheaves on $D$ to their component lying along the exceptional divisor $E$. In other words, the difference between these two functors is a sheaf on the proper transform $\widetilde{D}$ of $D$. The final pullback $j^*$ will kill this component, forcing the two functors to agree.

Thus, we have a well-defined functor $F\colon\cC\to \Coh(\frX)$.
To understand this functor,
we will need to apply
Orlov's result \cite{Orl92} that $p^*\Coh(X)$ and $q^*\Coh(H)\otimes\cO_E(1)$ generate the category $\Coh(\Bl_H(X))$ of coherent sheaves on the blowup, and in fact they form a semi-orthogonal decomposion:
\begin{equation}\label{eq:semiorth}
\Coh(\Bl_H(X)) = \langle \Coh(H)\otimes \cO_E(-1), \Coh(X)\rangle.
\end{equation}
Note that after pulling back to $\frX$, we can ignore the tensor
product with $\cO_E(-1)$ in \eqref{eq:semiorth}, since deleting $\widetilde{D}$ makes $E$ affine. This makes clear why $F$ is essentially surjective: the image of the functor $F$ contains the pullbacks to $\frX$ of both components in the semi-orthogonal decomposition~\eqref{eq:semiorth} and hence contains all objects in $j^*\Coh(\Bl_H(X))$.

To see why the functor $F$ is in fact an equivalence of categories, recall that the
category $\Coh(\frX)$ is a localization of $\Coh(\Bl_H(X))$ obtained by killing the identity morphism of any object in $\Coh\big(\widetilde{D}\big)$, and observe that this is equivalent to the localization which identifies any object in the second term of~\eqref{eq:semiorth} of the form~$k_*\cF$, for some~$\cF$ in $\Coh(D)$, with the object $(iq)^*\cF\otimes \cO_E(-1)$ in the first term.

Now note that the
category obtained from the semi-orthogonal decomposition \eqref{eq:semiorth} by identifying the images of $\Coh(X)$ and $\Coh(H)$ is equivalent to the category obtained from the ortho\-go\-nal sum $\Coh(H)\oplus \Coh(X)$ by the same identification;
but $\Coh(H)\oplus \Coh(X)$ is the~co\-pro\-duct of the categories $\Coh(H)$ and $\Coh(X)$, and the identification just described is the coequalizer of the respective maps from $D$ to these categories. In other words, this presentation exhibits $\Coh(\frX)$ as the pushout \eqref{eq:pushout}, and $F$ as the canonical equivalence $\cC\cong \Coh(\frX)$.
\end{proof}
\begin{thebibliography}{99}
\bibitem[24]{Orl92}
Orlov D.O., Projective bundles, monoidal transformations, and derived
 categories of coherent sheaves, \href{https://doi.org/10.1070/IM1993v041n01ABEH002182}{\textit{Russian Acad. Sci. Izv. Math.}}
 \textbf{41} (1993), 133--141.

\end{thebibliography}
\end{document}
